\begin{lemma}[The relative scheme of maximal rank-one spaces]
\label{lem:rank-one-fano-scheme}
Let \(B\) be any complex scheme and \(E,F\) vector bundles of
rank \(n\ge2\).  In \(\Gr_B(n,E\otimes F)\), let
\(Z(E,F)\) be the zero scheme of the restriction of the quadratic
rank-one equations to the universal subbundle \(\mathcal T\).
Equivalently it is the zero scheme of
\[
 \Sym^2\mathcal T\longrightarrow
              \bigwedge^2E\otimes\bigwedge^2F.
\]
The natural morphism
\begin{equation}\label{eq:fano-rank-one-scheme}
 \Pj_B(E)\ \sqcup\ \Pj_B(F)\longrightarrow Z(E,F),
 \qquad \ell\mapsto\ell\otimes F,
 \quad m\mapsto E\otimes m,
\end{equation}
is an isomorphism.  It commutes with arbitrary change of complex
base, including change to a nonreduced base.
\end{lemma}
\begin{proof}
We first work over \(\C\).  Each displayed map to the
Grassmannian is a closed immersion.  Its composition with the
Pluecker embedding is the degree-\(n\) Veronese embedding on
\(\ell^{\otimes n}\otimes\det F\), respectively
\(\det E\otimes m^{\otimes n}\), followed by its natural linear
inclusion into \(\bigwedge^n(E\otimes F)\).  This linear inclusion
is injective: polarization of the nonzero pure-power map gives the
corresponding irreducible symmetric-power summand.  The rank-one
equations vanish identically on both universal families, so the
closed immersions factor through \(Z(E,F)\).

Every geometric point is in one of the two images.  Indeed the sum
of two rank-one tensors is rank at most one only if their first
factors or their second factors are proportional.  Starting with two
independent tensors, this observation forces a linear space consisting
of rank-one tensors to fix the same factor throughout.  A space of
dimension \(n\) is then exactly one of the displayed spaces.
The two images are disjoint: a space fixing both factors has dimension
one, not \(n\).

We verify that these pointwise statements omit no infinitesimal
structure.  At \(\ell\otimes F\), choose a generator of \(\ell\)
and a complement in \(E\).  A Grassmannian tangent vector is a map
\(\phi:F\to(E/\ell)\otimes F\).  For each coordinate of
\(E/\ell\), write its component as \(A:F\to F\).
Linearizing the restricted minors gives
\[
                       u\wedge A(u)=0\quad\text{for every }u\in F.
\]
Taking the basis vectors shows that \(A\) is diagonal; taking their
pairwise sums makes all its diagonal entries equal.  Conversely a
scalar \(A\) satisfies all the linearized equations.  Thus the
Zariski tangent space to \(Z(E,F)\) is precisely
\(\Hom(\ell,E/\ell)\), of dimension \(n-1\).
The other image has the symmetric calculation.

At every closed point the closed immersed projective space gives
local dimension at least \(n-1\), whereas the tangent calculation
gives embedding dimension exactly \(n-1\).  The local ring is
therefore regular.  A finite-type complex scheme whose closed-point
local rings are regular is reduced: a nonempty support of its
nilradical would contain a closed point.  In particular there is no
nilpotent thickening or embedded associated structure at either
image.  The disjoint union of the two closed immersions is itself a
closed immersion, surjective on geometric points.  Its ideal in the
now reduced target vanishes at every closed point, so it is zero.
This proves the scheme isomorphism over \(\C\).

For general \(B\), trivialize \(E,F\) on a common open cover.
The Grassmannian and the zero scheme of the displayed universal
quadratic map are, on such an open, the base change of the constant
construction just proved.  Thus \eqref{eq:fano-rank-one-scheme} is
an isomorphism there even if the base has nilpotents.  The maps are
natural under changes of frame and hence glue.  This argument is
an identification by equations and base change, not an inference of
relative flatness from a fibrewise tangent calculation.
\end{proof}
